\begin{abstract}
Controllers transferred between traffic networks must rank actions under changes
in topology, demand and signal programs.  We introduce anchored Causal-Factored
Cross-City Model Transfer (CFCMT), a target-offline method that learns matched
counterfactual action contrasts.  From each restored state, a parent-restricted
rigid anchor and an antisymmetric all-pair model predict candidate-relative
costs; a seed-blocked target selector freezes their blend before deployment.
Target simulator labels make the setting calibration-light rather than
zero-shot.  Development on 18 SUMO networks from seven city groups reduced
leave-one-city-out normalized action regret by 9.84\% on an excluded seed, with
six of seven groups improving.  After method freeze, the untouched-seed regret
across two external networks was 0.269, compared with 0.292 for the
rigid anchor and 0.310 for a legacy near-target-only comparator.  In 56 matched
3,600-s rollouts, CFCMT reduced all-departed waiting by 6.27\% relative to a
same-information dense residual baseline.  A preregistered 64-seed comparison
then found a 3.04\% Jinan gain on all seed-level city contrasts but a 7.87\% Los
Angeles regression, rejecting joint two-city source benefit.  A post-V91
target-offline selector returned target-only control for Los
Angeles and admitted a Hangzhou source component for Jinan before fresh
rollouts.  On 56 further Jinan seeds, the admitted controller reduced waiting
by 4.53\% relative to a strict zero-source-row target model (95\% paired
interval, 4.28--4.79\%; 56/56 seeds improved).  Pressure policies
remained stronger.  These results support matched action contrast over dense
residual transfer and a protocol-specific V98 controller-pair gain in Jinan.
Because the V98 comparator is a lower-capacity zero-source-row target model,
that comparison does not isolate source-row contribution or establish universal
cross-city benefit or superiority to pressure control.
\end{abstract}
